\chapter{勘误与补充说明}

\label{app:errata}

正文直接采用正确写法；本附录记录实质修正及必要条件补充。原稿本来正确的结论不因列入此处而被视为错误；部分条目只是限制适用范围或明确记号。点原稿定位可返回正文；未删除原稿学习内容。

\section*{核对依据}

\begin{description}

\item[H1.] 汉森，\emph{计量经济学}，2010 年 1 月 10 日稿，第 2 章：条件均值、投影与解释。 \href{https://users.ssc.wisc.edu/~xshi28/Econometrics_Hansen.pdf}{大学网站公开的教材手稿}.

\item[H2.] 同一手稿，第 3.3 至 3.7 节及第 3.10 节：OLS 代数、残差回归与拟合优度。

\item[H3.] 同一手稿，第 4.3、4.4、4.11、4.12 节：条件无偏性、方差、共线性与遗漏变量。本笔记还明确推导了额外的有限样本限定条件。

\item[M.] 梅西大学， \href{https://r-resources.massey.ac.nz/161251/lectures/RegressionOrigin.html}{\emph{过原点回归}}: 中心化与非中心化拟合统计量的比较。

\item[计算.] 对原稿算式直接进行代数推导或数值运算；项目中包含可复现的数值核对。

\end{description}

外部材料仅用于核对数学与适用条件；本次排版模板仍是用户指定的旧 BMO 笔记。

\par\medskip\noindent\begin{minipage}{\linewidth}

\phantomsection\label{err:notation}

\textbf{C01. 总体、估计量与拟合值记号}\quad\textbf{原稿定位：}\hyperref[src:w2p4]{第 2 周，PDF 第 4 页}；\hyperref[src:w3p1]{第 3 周，PDF 第 1 页}。\quad\textbf{核对：}H1。\par

总体参数写作 $\beta_j$，样本估计写作 $\hat\beta_j$；总体误差 $u_i$ 与残差 $\hat u_i$ 分开。拟合值变化使用 $\Delta\hat y$，不将 $\hat\beta_j$ 直接写成总体参数。

\end{minipage}\par

\par\medskip\noindent\begin{minipage}{\linewidth}

\phantomsection\label{err:mean-normalisation}

\textbf{C02. 误差均值的归一化}\quad\textbf{原稿定位：}\hyperref[src:w2p1]{第 2 周，PDF 第 1 页}。\quad\textbf{核对：}H1。\par

令 $\varepsilon=u-\alpha$ 后，$\E[\varepsilon]=0$，而不是 $\varepsilon$ 处处等于零。截距变为 $\beta_0+\alpha$。该归一化本身不保证零条件均值。

\end{minipage}\par

\par\medskip\noindent\begin{minipage}{\linewidth}

\phantomsection\label{err:slope-correlation}

\textbf{C03. 斜率与相关系数的比例}\quad\textbf{原稿定位：}\hyperref[src:w2p2]{第 2 周，PDF 第 2 页}。\quad\textbf{核对：}H2。\par

样本简单 OLS 斜率为 $r_{xy}s_y/s_x$，不是 $r_{xy}s_x/s_y$；总体投影系数和样本系数不能混用。

\end{minipage}\par

\par\medskip\noindent\begin{minipage}{\linewidth}

\phantomsection\label{err:causality}

\textbf{C04. 相关与因果的边界}\quad\textbf{原稿定位：}\hyperref[src:w2p2]{第 2 周，PDF 第 2 页}；\hyperref[src:w2p4]{第 2 周，PDF 第 4 页}。\quad\textbf{核对：}H1。\par

保留“相关联”的措辞；控制变量、ZCM 或好的拟合度本身不能替代结构模型与可辩护的因果识别。截距在 $x=0$ 不可行时只是外推。

\end{minipage}\par

\par\medskip\noindent\begin{minipage}{\linewidth}

\phantomsection\label{err:wage-counterfactual}

\textbf{C05. 工资题中的单位与个体反事实}\quad\textbf{原稿定位：}\hyperref[src:w2p3]{第 2 周，PDF 第 3 页}。\quad\textbf{核对：}计算。\par

按原稿“千美元”复算：$5.58$、$-5.574$、$7.74$、$2.166$ 均使用同一单位。固定该人残差得到的是模型情景，并非保证的实际工资。未把看似异常的原题单位擅自改成美元。

\end{minipage}\par

\par\medskip\noindent\begin{minipage}{\linewidth}

\phantomsection\label{err:election}

\textbf{C06. 选举题的比例与金额}\quad\textbf{原稿定位：}\hyperref[src:w2p3]{第 2 周，PDF 第 3 页}。\quad\textbf{核对：}计算。\par

40\% 应输入 40，预测为 45.37\%；50\% 的预测阈值对应 支出占比 $>49.978448\ldots$。补齐在 B 支出 20 亿美元时 A 的支出阈值，区别预测过半与实际获胜。

\end{minipage}\par

\par\medskip\noindent\begin{minipage}{\linewidth}

\phantomsection\label{err:fwl-residual}

\textbf{C07. FWL 的单侧与双侧残差化}\quad\textbf{原稿定位：}\hyperref[src:w2p5]{第 2 周，PDF 第 5 页}；\hyperref[src:w4p5]{第 4 周，PDF 第 5 页}。\quad\textbf{核对：}H2。\par

只残差化自变量后回归原始 $y$，系数相同，但残差一般不等于原多元模型残差。对 $y$ 与该自变量同时排除其他回归量，才能同时得到相同系数与残差。

\end{minipage}\par

\par\medskip\noindent\begin{minipage}{\linewidth}

\phantomsection\label{err:ovb-levels}

\textbf{C08. 遗漏变量公式的三个层次}\quad\textbf{原稿定位：}\hyperref[src:w2p6]{第 2 周，PDF 第 6 页}；\hyperref[src:w4p8]{第 4 周，PDF 第 8 页}。\quad\textbf{核对：}H3。\par

保留样本恒等式 $\hat{\tilde\beta}_1=\hat\beta_1+\hat\beta_2\hat\delta_1$。给定整组设计时，偏差为 $\beta_2\hat\delta_1$；无条件有限样本偏差为 $\beta_2\E[\hat\delta_1]$。写成 $\beta_2\delta_1$ 还须有辅助斜率无偏等附加条件；总体投影差另行说明。

\end{minipage}\par

\par\medskip\noindent\begin{minipage}{\linewidth}

\phantomsection\label{err:unit-intercept}

\textbf{C09. 换单位时的截距与条件均值}\quad\textbf{原稿定位：}\hyperref[src:w2p7]{第 2 周，PDF 第 7 页}。\quad\textbf{核对：}计算。\par

总体截距的含义是 $\E[y\mid x=0]$，不是一般的 $\E[y]$；$y$ 乘以 $c$ 时截距也乘以 $c$。$x$ 乘以非零 $c$ 时斜率除以 $c$，截距不变。

\end{minipage}\par

\par\medskip\noindent\begin{minipage}{\linewidth}

\phantomsection\label{err:logs}

\textbf{C10. 有限变化与对数近似}\quad\textbf{原稿定位：}\hyperref[src:w2p8]{第 2 周，PDF 第 8 页}。\quad\textbf{核对：}计算。\par

将有限变化的近似与精确关系分开。水平值—对数模型为 $\beta_1\ln(x_1/x_0)$；对数—水平值模型的精确百分比为 $100(e^{\beta_1\Delta x}-1)$；对数—对数模型使用比例的幂。对数模型的系数仍保留原稿的局部百分比解释。

\end{minipage}\par

\par\medskip\noindent\begin{minipage}{\linewidth}

\phantomsection\label{err:sample-population}

\textbf{C11. 样本恒等式不是总体外生性的检验}\quad\textbf{原稿定位：}\hyperref[src:w3p1]{第 3 周，PDF 第 1 页}。\quad\textbf{核对：}H2。\par

残差和为零、与包含的回归量正交来自样本 OLS 一阶条件。不能从这些条件推出真实误差的 ZCM。总体均值关系用无帽参数，样本均值关系用样本系数。

\end{minipage}\par

\par\medskip\noindent\begin{minipage}{\linewidth}

\phantomsection\label{err:r2-degenerate}

\textbf{C12. 判定系数 的退化情形}\quad\textbf{原稿定位：}\hyperref[src:w3p2]{第 3 周，PDF 第 2 页}。\quad\textbf{核对：}H2。\par

补充 TSS 必须大于零。拟合值恒定时，普通 $R^2$ 可以为零，但相关系数因分母为零而未定义；此时不能直接用相关系数平方计算。

\end{minipage}\par

\par\medskip\noindent\begin{minipage}{\linewidth}

\phantomsection\label{err:r2-rss}

\textbf{C13. 高拟合度与 RSS 的比较}\quad\textbf{原稿定位：}\hyperref[src:w3p3]{第 3 周，PDF 第 3 页}。\quad\textbf{核对：}H2。\par

高 $R^2$ 表示较小的 RSS/TSS；不能跨不同样本或单位仅凭较高 $R^2$ 就认定绝对 RSS 更小。图中 0.94 和 0.64 均由新绘制的示意点实际计算。

\end{minipage}\par

\par\medskip\noindent\begin{minipage}{\linewidth}

\phantomsection\label{err:r2-monotonic}

\textbf{C14. 增加回归量时 判定系数 不减}\quad\textbf{原稿定位：}\hyperref[src:w3p4]{第 3 周，PDF 第 4 页}。\quad\textbf{核对：}H2。\par

原稿的“增加”改为“不下降，可以保持不变”，并注明必须是相同观测、相同因变量的嵌套 OLS 比较。不同缺失值样本不在该结论内。

\end{minipage}\par

\par\medskip\noindent\begin{minipage}{\linewidth}

\phantomsection\label{err:adjusted-meaning}

\textbf{C15. 调整判定系数 的提高条件}\quad\textbf{原稿定位：}\hyperref[src:w3p4]{第 3 周，PDF 第 4 页}。\quad\textbf{核对：}计算。\par

“变量有意义才提高”改写为残差均方必须下降；增加 $q$ 个回归量时，对应描述性的 偏 F 统计量 大于 1。这个阈值不是常规显著性检验，也不证明因果。

\end{minipage}\par

\par\medskip\noindent\begin{minipage}{\linewidth}

\phantomsection\label{err:demeaning}

\textbf{C16. 去均值与过原点回归}\quad\textbf{原稿定位：}\hyperref[src:w3p5]{第 3 周，PDF 第 5 页}；\hyperref[src:w3p6]{第 3 周，PDF 第 6 页}。\quad\textbf{核对：}H2。\par

是“去均值后无截距”的斜率等于“原始数据带截距”的斜率；不是任意无截距回归都等于带截距回归。原稿的两次推导均保留。

\end{minipage}\par

\par\medskip\noindent\begin{minipage}{\linewidth}

\phantomsection\label{err:rto-covariance}

\textbf{C17. 无截距时的正交与不相关}\quad\textbf{原稿定位：}\hyperref[src:w3p7]{第 3 周，PDF 第 7 页}。\quad\textbf{核对：}H2。\par

$\sum_i x_i\check u_i=0$ 仍成立，但 $\Cov_n(x,\check u)=-\bar x\bar{\check u}$ 一般不为零。原稿将两者等同的说法已改正。

\end{minipage}\par

\par\medskip\noindent\begin{minipage}{\linewidth}

\phantomsection\label{err:rto-cross-term}

\textbf{C18. 平方和分解的交叉项}\quad\textbf{原稿定位：}\hyperref[src:w3p7]{第 3 周，PDF 第 7 页}。\quad\textbf{核对：}H2。\par

交叉项应为 $2\sum_i\check u_i(\hat y_i-\bar y)$，其中是拟合值 $\hat y_i$。无截距时它不一定消失，因此不能一般地写 TSS=ESS+RSS。

\end{minipage}\par

\par\medskip\noindent\begin{minipage}{\linewidth}

\phantomsection\label{err:rto-r2}

\textbf{C19. 有截距与无截距 判定系数 的基准}\quad\textbf{原稿定位：}\hyperref[src:w3p7]{第 3 周，PDF 第 7 页}。\quad\textbf{核对：}M。\par

区分以样本均值为基准的 中心化判定系数 与以零为基准的 非中心化判定系数。前者在过原点回归中可以为负；两种统计量不能当作同一尺度直接比较。

\end{minipage}\par

\par\medskip\noindent\begin{minipage}{\linewidth}

\phantomsection\label{err:estimator-estimate}

\textbf{C20. 估计量与某次估计值}\quad\textbf{原稿定位：}\hyperref[src:w4p1]{第 4 周，PDF 第 1 页}。\quad\textbf{核对：}H3。\par

无偏性不要求某一次估计等于真值，也不声称每一次估计必定不等于真值。它是对重复抽样期望的限制，并需要相关期望存在。

\end{minipage}\par

\par\medskip\noindent\begin{minipage}{\linewidth}

\phantomsection\label{err:conditioning}

\textbf{C21. 无偏性证明的条件期望}\quad\textbf{原稿定位：}\hyperref[src:w4p2]{第 4 周，PDF 第 2 页}；\hyperref[src:w4p5]{第 4 周，PDF 第 5 页}。\quad\textbf{核对：}H3。\par

先给定整组样本设计，再将样本权重与分母作为常量提出条件期望。不能仅使用单个 $x_i$ 的条件，或在无条件期望中把随机分母当常数。注明样本满秩和期望存在。

\end{minipage}\par

\par\medskip\noindent\begin{minipage}{\linewidth}

\phantomsection\label{err:intercept-sign}

\textbf{C22. 截距证明的正负号}\quad\textbf{原稿定位：}\hyperref[src:w4p3]{第 4 周，PDF 第 3 页}。\quad\textbf{核对：}计算。\par

从总体模型取平均得到 $\bar y=\beta_0+\beta_1\bar x+\bar u$；因此截距展开式中的 $\bar u$ 应为加号。原稿中间一行的减号与其余推导不一致。

\end{minipage}\par

\par\medskip\noindent\begin{minipage}{\linewidth}

\phantomsection\label{err:violations}

\textbf{C23. 违背抽样假设的后果}\quad\textbf{原稿定位：}\hyperref[src:w4p3]{第 4 周，PDF 第 3 页}。\quad\textbf{核对：}H3。\par

保留时间序列相关、非代表性样本和溢出效应的例子；将它们与“必然产生 OLS 偏差”区分。需另查条件均值，不能仅凭违反独立同分布假设 就推出所有系数有偏。

\end{minipage}\par

\par\medskip\noindent\begin{minipage}{\linewidth}

\phantomsection\label{err:zcm-covariance}

\textbf{C24. ZCM 强于零协方差}\quad\textbf{原稿定位：}\hyperref[src:w4p3]{第 4 周，PDF 第 3 页}；\hyperref[src:w4p5]{第 4 周，PDF 第 5 页}。\quad\textbf{核对：}H1。\par

$\E[u\mid x]=0$ 推出零均值、零协方差；反方向不成立。总体误差不相关不足以保证有限样本无偏；与样本残差正交也不构成外生性的证据。

\end{minipage}\par

\par\medskip\noindent\begin{minipage}{\linewidth}

\phantomsection\label{err:collinearity}

\textbf{C25. 完全共线性与平方项}\quad\textbf{原稿定位：}\hyperref[src:w4p4]{第 4 周，PDF 第 4 页}。\quad\textbf{核对：}H3。\par

保留常数、倍数、仿射关系、总支出等例子。对数变量及其平方一般不自动完全共线，但特殊样本仍可能不满秩；“样本满秩所以总体一定如何”的措辞改为直接检查样本设计。

\end{minipage}\par

\par\medskip\noindent\begin{minipage}{\linewidth}

\phantomsection\label{err:unbiased-log}

\textbf{C26. 9.2\% 与无偏估计程序}\quad\textbf{原稿定位：}\hyperref[src:w4p6]{第 4 周，PDF 第 6 页}。\quad\textbf{核对：}H3。\par

改正 $\log(\widehat{wage})$ 与 $\widehat{\log(wage)}$ 的混用。0.092 是一个实现值；OLS 系数估计量的无偏不保证指数变换无偏。9.2\% 是近似，精确的固定误差变换约为 9.64\%。

\end{minipage}\par

\par\medskip\noindent\begin{minipage}{\linewidth}

\phantomsection\label{err:irrelevant-exogeneity}

\textbf{C27. 加入无关变量需要扩充后的外生性}\quad\textbf{原稿定位：}\hyperref[src:w4p7]{第 4 周，PDF 第 7 页}。\quad\textbf{核对：}H3。\par

不能从 $\E[u\mid x_1,x_2]=0$ 自动推出加入 $x_3$ 后条件均值仍为零。仅在扩充后满秩且继续满足 ZCM 等条件时，原有系数及零系数的 OLS 估计量才保持无偏。

\end{minipage}\par

\par\medskip\noindent\begin{minipage}{\linewidth}

\phantomsection\label{err:irrelevant-variance}

\textbf{C28. 无关控制变量与精度}\quad\textbf{原稿定位：}\hyperref[src:w4p7]{第 4 周，PDF 第 7 页}。\quad\textbf{核对：}H3。\par

将“估计方差必然上升”限定为同方差固定设计下，加入无关控制使相关真实条件方差弱增；不保证严格增加，也不保证每次实现的估计标准误增加。

\end{minipage}\par

\par\medskip\noindent\begin{minipage}{\linewidth}

\phantomsection\label{err:omitting-relevant}

\textbf{C29. 遗漏相关变量的条件}\quad\textbf{原稿定位：}\hyperref[src:w4p7]{第 4 周，PDF 第 7 页}。\quad\textbf{核对：}H3。\par

小节标题按内容改为 遗漏相关变量；遗漏变量系数非零并不足以保证目标斜率改变，还要看它与保留回归量的关系。误差、短回归残差、长回归残差分开。

\end{minipage}\par

\par\medskip\noindent\begin{minipage}{\linewidth}

\phantomsection\label{err:toward-zero}

\textbf{C30. 向零偏不能只看上偏或下偏}\quad\textbf{原稿定位：}\hyperref[src:w4p8]{第 4 周，PDF 第 8 页}。\quad\textbf{核对：}计算。\par

保留正系数下偏、负系数上偏的方向解释，但增加“不越过零”的限制。发生符号反转时不能仅按方向就称为通常意义的向零偏。

\end{minipage}\par
